\documentclass[11pt,a4paper,twocolumn]{article}
\usepackage[margin=21mm,headheight=14pt,includeheadfoot]{geometry}
\usepackage{amsmath,booktabs,enumitem,fancyhdr}
\usepackage[hidelinks]{hyperref}

\setlength{\columnsep}{6mm}
\setlength{\parindent}{0pt}
\setlength{\parskip}{3pt}
\setlist[enumerate]{label=\textbf{(\alph*)},leftmargin=*,itemsep=3pt,topsep=2pt,parsep=0pt}
\pagestyle{fancy}
\fancyhf{}
\fancyhead[L]{\small COMP2040 --- Operating Systems}
\fancyhead[R]{\small Checkpoint 1}
\fancyfoot[C]{\small\thepage}
\newcommand{\code}[1]{\texttt{#1}}

% Results are taken from the merged report and checked-in validation notes.
% No simulator or LaTeX builds were run while preparing this revision.
\begin{document}
\twocolumn[{%
\begin{center}
  {\Large\bfseries COMP2040: Operating Systems\par}
  \vspace{3pt}
  {\large Checkpoint 1: Progress Report --- Group 10\par}
  \vspace{4pt}
  Nguyen Vu Nhat Minh (V202401314) \quad Do Chi Hai (V202401459)\\
  {\small College of Engineering and Computer Science, VinUniversity}\\
  {\small Repository: \url{https://github.com/TheNameisSea/COMP2040_Project}}
\end{center}
\vspace{2pt}
}]

We have implemented the four required schedulers and resource deadlock
detection. The following explains the design, the checks we used, and
the remaining limitations.

\section{Schedulers}

SJF, SRTF and PRIO each have their own \code{queue\_t}: an array of
pointers to process control blocks, with space for 64 processes. A READY
process joins the tail. \code{q\_pop\_min} picks the smallest key; ties
follow queue order, so the process that became ready first wins.

\begin{enumerate}
\item \textbf{SJF.} The key is the remaining length of the current CPU
burst, not the process's total future work. SJF never preempts. A process
keeps the CPU until it blocks or finishes; consecutive CPU operations
continue without a new dispatch.

\item \textbf{SRTF.} It uses the same burst-length key in a separate
queue. Before each tick, it preempts only if a ready process has
\emph{strictly less} time left than the runner. Equal lengths do not
preempt. Arrivals and I/O completions are handled first, so they can
change the decision in that tick.

\item \textbf{PRIO.} The key is \code{priority}: a smaller number means
higher priority. A higher priority preempts immediately. Equal
priorities take turns after \(q\) CPU ticks. On expiry, the runner joins
the tail only if an equal priority is ready; lower priorities cannot
force a rotation. The PCB's \code{quantum\_used} counter advances on CPU
ticks and resets on every dispatch.

\item \textbf{MLFQ.} Three FIFO queues represent levels 0, 1 and 2,
with quanta \(q\), \(2q\) and \(4q\). We pick the head of the highest
non-empty queue. A full quantum moves a process down one level; at
level 2 it starts another quantum. Blocking earlier keeps its level.
A higher-level ready process preempts immediately. After expiry, the
runner gives way to a process at its new level or higher.

At ticks 20, 40, and so on, a boost puts every live process, including
blocked and running ones, at level 0. Lower queues move to level 0 in
FIFO order. Both \code{pick} and \code{should\_preempt} call the boost
helper, so it also runs when the CPU is idle.
\end{enumerate}

Minimum selection and preemption scans take \(O(n)\), which is adequate
for 64 processes. The simulator knows exact burst lengths; a real OS
could estimate them by exponential averaging.

\section{Deadlock Detection}

We use the multi-instance algorithm from textbook Section~8.6.2.
\code{Work} starts as the available resources. Processes holding nothing,
and terminated processes, are marked finished. We repeatedly find an
unfinished process whose outstanding Request fits \code{Work}, assume
it finishes, and add its Allocation back to \code{Work}.

Processes left unfinished are deadlocked. Their PIDs are printed in
ascending order, for example \code{DEADLOCK t=15 pids: 1 2}. We remember
the last reported set and print again only when it changes. Detection
takes \(O(n^2m)\) for \(n\) processes and \(m\) resource types.

\section{Validation}

We worked out Gantt charts and per-process times for
\code{workloads/arrivals.txt} before comparing with the simulator.
Paging was off; PRIO and MLFQ used \(q=2\). Waiting counts READY ticks,
response is first dispatch minus arrival, and turnaround is finish
minus arrival.

\begin{center}
{\footnotesize
\setlength{\tabcolsep}{3pt}
\begin{tabular}{@{}lrrr@{}}
\toprule
 & \multicolumn{3}{c}{Average ticks: hand / \code{ossim}} \\
\cmidrule(l){2-4}
Policy & Waiting & Response & Turnaround \\
\midrule
SJF  & 10.00/10.00 & 9.00/9.00 & 17.83/17.83 \\
SRTF & 9.83/9.83   & 9.00/9.00 & 17.67/17.67 \\
PRIO & 12.83/12.83 & 8.33/8.33 & 20.67/20.67 \\
MLFQ & 15.33/15.33 & 0.00/0.00 & 23.17/23.17 \\
\bottomrule
\end{tabular}}
\end{center}

All three averages and context-switch counts are checked in
\code{tests/run.sh}. Full traces are in \code{docs/checks/}.
The MLFQ trace includes the boost at tick 20 and gives 19 switches,
counting the initial dispatch.

Additional workloads check FIFO ties, consecutive bursts, PRIO quantum
resets, and MLFQ demotion and boosting. Producer/consumer checks pass
under RR and all four new policies. Mixed paging workloads also finish,
including runs with one frame.

For \code{deadlock.txt}, B blocks at the end of tick 11 and A at the end
of tick 12. With checks every five ticks, the first report is at tick 15.
In the two-instance case, A is blocked at tick 4 but B currently requests
zero: letting B finish hypothetically frees an instance for A. Both
request another instance by tick 6, when detection reports the deadlock.

\textbf{Terminal excerpts.}
The recorded GCC validation used
\code{-std=c99 -Wall -Wextra -Werror}. The supplied RR baseline on
\code{basic.txt}, with quantum 4 and paging off, produced:

\begingroup
\footnotesize\ttfamily\raggedright
\$ ./ossim workloads/basic.txt --sched rr --quantum 4 --no-paging --gantt\par
finished=3/3 context\_switches=4\par
avg\_turnaround=15.67 avg\_waiting=5.67 avg\_response=3.67\par
\$ make test\par
...\par
65 passed, 4 skipped (not implemented yet), 0 failed\par
\endgroup

The hand averages agree: waiting \((6+4+7)/3=5.67\), response
\((0+4+7)/3=3.67\), and turnaround \((30+7+10)/3=15.67\) ticks.
The \code{arrivals.txt} table above separately checks our four new
schedulers.

Both test helpers reject command errors. In particular, the negative
check cannot mistake an error message for a successful run.

\section{Limits and Rule Choices}

All Checkpoint 1 algorithms are implemented; the four SKIPs belong to
Checkpoint 2. PRIO has no aging, so low priorities can starve. Detection
covers resources, not semaphore-only or all mixed deadlocks. Recovery is
optional and has not been added. Clang has not been verified. The suite
also lacks a dedicated FCFS producer/consumer completion check.

The handout leaves some timing details open. We chose fresh quanta on
dispatch, including after higher-priority interruption. MLFQ demotes on
a full quantum even if the process then blocks; a boost resets the
runner's quantum and appends level 1, then level 2, behind existing
level-0 waiters.

\section{Contributions}

\textbf{Nguyen Vu Nhat Minh:} SJF/SRTF/PRIO, their hand checks and
regression workloads, stricter test helpers, WSL line-ending fixes, and
resource-deadlock review.

\textbf{Do Chi Hai:} MLFQ and deadlock detection, their hand checks,
workloads and tests, plus the team guide and decision log.

AI assistants helped with implementation, hand checks and debugging.

\end{document}
